\begin{tomtat}
\begin{itemize}
  \item Hiệu năng dự đoán không tự trả lời các câu hỏi về cơ chế, mục tiêu còn
  thiếu trong đặc tả, hay cách con người sẽ dùng mô hình.
  \item \tn{Khả năng diễn giải}{interpretability} là một mục tiêu gắn với
  người nhận và công việc; \tn{minh bạch}{transparency} và \tn{giải thích hậu
  nghiệm}{post-hoc explanation} có nguồn gốc khác nhau.
  \item Đánh giá trong ứng dụng, đánh giá với người dùng và đánh giá theo chức
  năng dùng các loại bằng chứng khác nhau.
  \item \tn{Tính hợp lý}{plausibility} với người đọc không chứng minh
  \tn{độ trung thực}{faithfulness} với cơ chế mô hình.
  \item Một chỉ số độ trung thực là dụng cụ đo và cần được kiểm chứng như một
  dụng cụ đo.
\end{itemize}
\end{tomtat}

\cauhoi
\begin{cauhoids}
  \item Nêu một tình huống trong đó mô hình đạt hiệu năng tốt nhưng người dùng
  vẫn cần một lời giải thích. Xác định điều gì còn thiếu trong đặc tả bài toán.
  \item Chọn một lời giải thích hậu nghiệm quen thuộc. Theo
  định nghĩa~\ref{def:ch01-do-trung-thuc}, hãy nêu hành vi và loại can thiệp
  cần xác định trước khi có thể hỏi nó có trung thực hay không.
  \item Một đánh giá theo chức năng cho thấy mô hình có ít quy tắc hơn mô hình
  khác. Kết quả đó hỗ trợ được kết luận nào, và không hỗ trợ được kết luận nào
  về việc người dùng có hiểu mô hình hay không?
\end{cauhoids}
